% =============================================================================
% CAPÍTULO 5: PROCESSOS TÉCNICOS
% Alinhado à ISO/IEC/IEEE 12207:2021 (6.4) e MPS.BR
% Fonte: manual-processos/manual/capitulos/05-processos-tecnicos.md
% =============================================================================

\aberturacapitulo{Capítulo 5}{Processos Técnicos}

\secao{Processo de definição de requisitos}

A Engenharia de Requisitos da \termochave{CineSoft Solutions} segue o processo formal da \termochave{ISO/IEC/IEEE 12207:2021 (6.4.2 e 6.4.3)}, garantindo que as necessidades dos clientes e regulamentos legais (como a LGPD) sejam traduzidas em requisitos verificáveis.

\subsecao{Catálogo Formal de Requisitos Funcionais e Não Funcionais}

\begin{center}
\small
\begin{tabularx}{\textwidth}{l p{3.2cm} c X}
\toprule
\textbf{ID} & \textbf{Nome do Requisito} & \textbf{Tipo} & \textbf{Descrição e Regra Central} \\
\midrule
RF-01 & Reserva Atômica de Poltrona & Funcional & Bloqueio temporário (10 min) prevenindo duplicidade. \\
RF-02 & Emissão e Baixa de Ingressos & Funcional & Geração de QR Code criptográfico para catraca/guichê. \\
RF-03 & Venda no Balcão (PDV Rápido) & Funcional & Operação simplificada para atendente em menos de 30s. \\
RNF-01 & Tempo de Resposta em Checkout & Não Func. & P95 inferior a 300ms sob carga de 50 req/s. \\
RNF-02 & Conformidade de Segurança & Não Func. & Sanitização Zod e isolamento total de tenants. \\
\bottomrule
\end{tabularx}
\end{center}

\subsecao{Matriz de Rastreabilidade e Critérios de Aceitação}

% =============================================================================
% ARTEFATO: Matriz de Rastreabilidade Bidirecional Completa (CineDoom)
% Capitulo 5 - secao 5.1.2 (Rastreabilidade de requisitos)
%
% Este arquivo contem a versao COMPLETA da matriz, incluindo a coluna de
% codigo-fonte. No corpo do manual (5.1.2) a matriz aparece em versao compacta
% (4 colunas) para preservar a legibilidade, remeter a este artefato.
%
% Requer os pacotes: booktabs, colortbl, tabularx, array
% =============================================================================

\begin{center}
\small
\setlength{\tabcolsep}{4pt}
\renewcommand{\arraystretch}{1.15}
\begin{tabularx}{\textwidth}{>{\raggedright\arraybackslash}m{2.6cm} >{\raggedright\arraybackslash}m{1.7cm} >{\raggedright\arraybackslash}m{1.3cm} >{\raggedright\arraybackslash}m{4.0cm} Y}
\arrayrulecolor{GithubTableBorder}\toprule
\rowcolor{BrandTableHeader}
\textbf{Necessidade / Stakeholder} & \textbf{Requisito(s)} & \textbf{Aceite} & \textbf{Codigo-Fonte (cinedoom-app/)} & \textbf{Suíte de Teste (tests/)} \\
\midrule
\arrayrulecolor{GithubGrayRule}
\texttt{STK-02} / \texttt{NEC-02} (Controle de Acesso e Equipe) & \texttt{RF-01}, \texttt{RF-02} & \texttt{CA-05} & \texttt{src/modules/auth/}, \texttt{src/middleware.ts}, \texttt{src/app/api/backoffice/usuarios/} & \texttt{auth-rbac.test.ts} (→ \texttt{RF-01}, \texttt{CA-05}); \texttt{users-management.test.ts} (→ \texttt{RF-02}, \texttt{CA-05}) \\
\midrule
\texttt{STK-02} / \texttt{NEC-02} (Gestao de Acervo e Salas) & \texttt{RF-03} a \texttt{RF-05}, \texttt{RES-05} & \texttt{CA-08} & \texttt{src/modules/backoffice/} (film-service, room-service, session-service), \texttt{src/app/api/} & \texttt{catalog-filmes.test.ts} (→ \texttt{RF-03}); \texttt{rooms-layout.test.ts} (→ \texttt{RF-04}); \texttt{sessions-prepopulation.test.ts} (→ \texttt{RF-05}) \\
\midrule
\texttt{STK-03}, \texttt{STK-04} / \texttt{NEC-03}, \texttt{NEC-06} (Mapa Reativo e Guiche) & \texttt{RF-06}, \texttt{RNF-01}, \texttt{RNF-04} & \texttt{CA-02} & \texttt{src/modules/bilheteria/components/seating-grid.tsx}, \texttt{src/modules/bilheteria/hooks/use-session-seats.ts}, \texttt{src/app/api/sessoes/[id]/assentos/} & \texttt{seating-grid.test.ts} (→ \texttt{RF-06}, \texttt{RNF-04}, \texttt{CA-02}); \texttt{session-seats-polling.test.ts} (→ \texttt{RF-06}, \texttt{RNF-01}, \texttt{CA-02}) \\
\midrule
\texttt{STK-01}, \texttt{STK-03} / \texttt{NEC-01}, \texttt{NEC-03} (Venda Atomica ACID) & \texttt{RF-07}, \texttt{RF-08}, \texttt{RNF-02}, \texttt{RES-01}, \texttt{RES-02}, \texttt{RES-04} & \texttt{CA-01}, \texttt{CA-15} & \texttt{src/modules/bilheteria/sale-service.ts}, \texttt{src/modules/bilheteria/counter-service.ts}, \texttt{src/app/api/vendas/}, \texttt{src/app/api/ingressos/}, \texttt{prisma/schema.prisma} & \texttt{venda-atomica.test.ts} (→ \texttt{RF-07}, \texttt{RNF-02}, \texttt{CA-01}); \texttt{concurrency-acid-50-threads.test.ts} (→ \texttt{RF-07}, \texttt{RNF-02}, \texttt{CA-01}); \texttt{cancelamento-auditavel.test.ts} (→ \texttt{RF-08}, \texttt{CA-15}) \\
\midrule
\texttt{STK-03}, \texttt{STK-05} / \texttt{NEC-04}, \texttt{NEC-08} (Cancelamento Imutavel) & \texttt{RF-09} & \texttt{CA-03} & \texttt{src/modules/bilheteria/cancellation-service.ts}, \texttt{src/app/api/ingressos/[id]/cancelar/}, \texttt{prisma/triggers.sql}, \texttt{prisma/schema.prisma} & \texttt{cancelamento-auditavel.test.ts} (→ \texttt{RF-09}, \texttt{CA-03}); \texttt{database-schema.test.ts} (→ \texttt{RF-09}, \texttt{CA-03}) \\
\midrule
\texttt{STK-05} / \texttt{NEC-04} (Privacidade e LGPD Art. 18) & \texttt{RF-10}, \texttt{RNF-03} & \texttt{CA-06} & \texttt{src/modules/crypto/crypto-service.ts}, \texttt{src/modules/compliance/}, \texttt{src/modules/compradores/buyer-service.ts} & \texttt{cryptoservice.test.ts} (→ \texttt{RF-10}, \texttt{RNF-03}, \texttt{CA-06}); \texttt{compliance-anonymization.test.ts} (→ \texttt{RF-10}, \texttt{CA-06}) \\
\midrule
\texttt{STK-01} / \texttt{NEC-01} (Isolamento Multi-Tenant SaaS) & \texttt{RNF-03}, \texttt{RNF-06} & \texttt{CA-04} & \texttt{src/modules/tenancy/prisma-tenancy-extension.ts}, \texttt{src/modules/tenancy/tenancy-context.ts}, \texttt{src/shared/db/prisma.ts} & \texttt{tenancy-isolation.test.ts} (→ \texttt{RNF-03}, \texttt{CA-04}); \texttt{security-cross-tenant-suite.test.ts} (→ \texttt{RNF-03}, \texttt{CA-04}); \texttt{modules.test.ts}, \texttt{ci-workflow.test.ts} (→ \texttt{RNF-06}) \\
\midrule
\texttt{STK-02}, \texttt{STK-05} / \texttt{NEC-02}, \texttt{NEC-04} (Painel de Ocupacao e Logs) & \texttt{RF-11} & \texttt{CA-07} & \texttt{src/app/api/ocupacao/painel/}, \texttt{src/modules/ocupacao/ocupacao-service.ts}, \texttt{src/shared/logger/index.ts} & \texttt{observability-logger.test.ts} (→ \texttt{RF-11}, \texttt{CA-07}); \texttt{security-cross-tenant-suite.test.ts} (→ \texttt{RF-11}, \texttt{CA-07}) \\
\midrule
\texttt{STK-04} / \texttt{NEC-05}, \texttt{NEC-06} (Portal e Reserva Web) & \texttt{RF-12} a \texttt{RF-14}, \texttt{RNF-04}, \texttt{RNF-05}, \texttt{RES-01}, \texttt{RES-03}, \texttt{RES-06} & \texttt{CA-09}, \texttt{CA-10}, \texttt{CA-15} & \texttt{src/modules/portal/portal-service.ts}, \texttt{src/app/portal/}, \texttt{src/app/api/portal/}, \texttt{src/app/layout.tsx} & \texttt{reserva-atomica.test.ts} (→ \texttt{RF-13}, \texttt{CA-09}); \texttt{cryptoservice.test.ts} (→ \texttt{RF-14}, \texttt{CA-10}); \texttt{auth-rbac.test.ts} (→ \texttt{RF-14}, \texttt{CA-10}); \texttt{seating-grid.test.ts} (→ \texttt{RNF-04}) \\
\midrule
\texttt{STK-04}, \texttt{STK-03} / \texttt{NEC-07}, \texttt{NEC-08} (Voucher, Cancelamento Web, Baixa e JIT) & \texttt{RF-15} a \texttt{RF-18}, \texttt{RES-01}, \texttt{RES-02} & \texttt{CA-11} a \texttt{CA-14} & \texttt{src/modules/notificacoes/notification-service.ts}, \texttt{src/app/api/portal/reservas/}, \texttt{src/app/api/bilheteria/reservas/}, \texttt{src/modules/bilheteria/counter-service.ts} & \texttt{reserva-atomica.test.ts} (→ \texttt{RF-15} voucher, \texttt{RF-16} consulta, \texttt{RF-18}, \texttt{CA-11}, \texttt{CA-14}); \textbf{sem automacao:} \texttt{RF-15} (e-mail/\texttt{email\_sent\_at}), \texttt{RF-16}/\texttt{CA-12} (cancelamento web), \texttt{RF-17}/\texttt{CA-13} (baixa presencial) \\
\arrayrulecolor{GithubTableBorder}\bottomrule
\end{tabularx}
\end{center}


\begin{center}
\small
\begin{tabularx}{\textwidth}{l p{3.2cm} X}
\toprule
\textbf{ID} & \textbf{Requisito Funcional} & \textbf{Critério de Aceitação Formal (Dado / Quando / Então)} \\
\midrule
CA-01 & Bloqueio Concorrente & \textbf{Dado} um assento livre, \textbf{quando} dois clientes requisitarem bloqueio no mesmo milissegundo, \textbf{então} apenas uma transação obtém sucesso (\texttt{200}) e a outra recebe conflito (\texttt{409}). \\
CA-02 & Expiração Automática & \textbf{Dado} um assento bloqueado há mais de 10 min sem checkout, \textbf{quando} ocorrer nova consulta, \textbf{então} o status reverte automaticamente para disponível. \\
CA-03 & Venda no Balcão PDV & \textbf{Dado} o operador no guichê, \textbf{quando} selecionar assento vago e confirmar pagamento em dinheiro/Pix, \textbf{então} o ingresso é emitido e impresso em $<3$ segundos. \\
\bottomrule
\end{tabularx}
\end{center}

\secao{Processo de arquitetura e design do sistema}

A arquitetura do CineDoom adota o padrão de \destaqueazul{Monólito Modular com Clean Architecture}, desacoplando regras de negócio corporativas de detalhes de infraestrutura e bibliotecas externas.

\subsecao{Diagrama de Contêineres (C4 Model Nível 2)}

O sistema CineDoom é modelado no C4 Model detalhando o relacionamento entre os contêineres:

\figuracomfonte{capitulo-05/diagramas/c4-containers.pdf}{0.96\textwidth}{Diagrama C4 Nível 2: Contêineres da Plataforma CineDoom}{CineSoft Solutions (2026)}{fig:c4_containers}

O sistema é composto por:
\begin{listapontos}
  \itemponto{Single Page / Server Component App}{Interface Next.js 16 servida na borda pela Vercel, proporcionando renderização híbrida rápida para mobile e guichê.}
  \itemponto{API Transacional de Negócios}{Camada de serviços escrita em TypeScript que orquestra a máquina de estados de reserva e venda.}
  \itemponto{Repositório Relacional Serverless}{Banco de dados Neon PostgreSQL conectado via Prisma ORM com suporte a transações ACID concorrentes.}
  \itemponto{Provedor de E-mail Externo}{Serviço Resend integrado via HTTPS REST para despacho assíncrono de notificações transacionais.}
\end{listapontos}

\subsecao{Modelo de Dados (DER) e Modelo de Domínio}

O modelo relacional e as entidades de domínio são formalizados nas figuras abaixo:

\figuracomfonte{capitulo-05/diagramas/diagrama-er-banco.pdf}{0.96\textwidth}{Diagrama Entidade-Relacionamento (DER) do Banco PostgreSQL}{CineSoft Solutions (2026)}{fig:der_banco}

\figuracomfonte{capitulo-05/diagramas/diagrama-classes-dominio.pdf}{0.96\textwidth}{Diagrama de Classes do Domínio de Negócios CineDoom}{CineSoft Solutions (2026)}{fig:classes_dominio}

\subsecao{Diagramas de Sequência, Casos de Uso e Máquina de Estados}

A máquina de estados, o mapeamento de casos de uso e o fluxo de reserva atômica são formalizados nas figuras abaixo:

\figuracomfonte{capitulo-05/diagramas/diagrama-casos-de-uso.pdf}{0.96\textwidth}{Diagrama de Casos de Uso UML da Plataforma CineDoom}{CineSoft Solutions (2026)}{fig:casos_uso}

\figuracomfonte{capitulo-05/diagramas/diagrama-estados-assento.pdf}{0.96\textwidth}{Diagrama de Máquina de Estados do Ciclo de Vida da Poltrona}{CineSoft Solutions (2026)}{fig:estados_assento}

\figuracomfonte{capitulo-05/diagramas/diagrama-sequencia-reserva.pdf}{0.96\textwidth}{Diagrama de Sequência UML: Reserva Web e Emissão de Voucher}{CineSoft Solutions (2026)}{fig:sequencia_reserva}

\secao{Processo de implementação e integração de software}

A implementação do CineDoom orienta-se pela manutenibilidade, cobertura de testes e aderência aos padrões de código:
\begin{listapontos}
  \itemponto{Padrões de Código}{Padronização rigorosa com ESLint e Prettier aplicada automaticamente via Git Hooks (Husky).}
  \itemponto{Migrations de Banco}{Todas as alterações de schema são versionadas no Prisma e aplicadas via pipeline automatizada no deploy de homologação.}
  \itemponto{Tratamento de Exceções}{Mapeamento de erros de concorrência com respostas padronizadas em JSON para o cliente web.}
\end{listapontos}

\figuracomfonte{capitulo-05/prints/print-testes-unitarios-vitest.png}{0.94\textwidth}{Terminal de Execução da Suíte Vitest (748 Testes Aprovados)}{CineSoft Solutions (2026)}{fig:testes_vitest}

\secao{Processos de verificação, validação e transição (V\&V)}

Em conformidade com a ISO 12207 (6.4.9 e 6.4.10), a CineSoft estabelece distinção clara entre Verificação e Validação:

\begin{center}
\small
\begin{tabularx}{\textwidth}{>{\fontSemiBold\color{BrandNavy}\arraybackslash}m{3.2cm} Y Y}
\arrayrulecolor{GithubTableBorder}\toprule
\rowcolor{BrandTableHeader}
\textbf{Dimensão} & \textbf{Verificação (Construir Certo)} & \textbf{Validação (Construir a Coisa Certa)} \\
\midrule
\arrayrulecolor{BrandGrayRule}
\textbf{Foco Central} & O software atende às especificações técnicas e regras de arquitetura? & O software resolve a necessidade real da bilheteria e do cliente final? \\
\midrule
\textbf{Técnicas} & Testes unitários Vitest, linting ESLint, checagem estática de tipos TypeScript. & Testes de usabilidade com atendentes reais, homologação UAT e pilotos. \\
\midrule
\textbf{Critério de Sucesso} & 748 testes unitários passando e zero erros de tipo no build. & Atendente fecha venda em menos de 30s e índice de erro nulo em guichê. \\
\arrayrulecolor{GithubTableBorder}\bottomrule
\end{tabularx}
\end{center}

\secao{Processos de operação, manutenção e descontinuação}

A operação do CineDoom inclui monitoramento de saúde do banco de dados, gestão de backups diários e procedimentos de rollback imediato em caso de incidentes.

\verificacaoconformidade{Ponto de Controle — Processos Técnicos}{%
  \begin{enumerate}[label=\textbf{\alph*)}, leftmargin=0.8cm, itemsep=3pt]
    \item Como a arquitetura modular do CineDoom garante que falhas no envio de e-mails pelo Resend não interrompam a venda física de ingressos na bilheteria?
    \item Quais são os critérios de entrada e saída para que uma nova versão do sistema avance do ambiente de homologação para produção?
  \end{enumerate}
}

\notatecnica{%
  \itemcaixa{O isolamento multi-tenant do CineDoom utiliza cláusula obrigatória de \codigoinline{tenantId} em todas as consultas Prisma, garantindo separação lógica absoluta entre os dados de diferentes redes exibidoras.}
}

\saibamaisnormativo{%
  \itemsaibamais{Consulte as cláusulas 6.4.4 (\textit{Architecture Definition}), 6.4.9 (\textit{Verification}) e 6.4.10 (\textit{Validation}) da norma \linkcaixa{https://www.iso.org/standard/63711.html}{ISO/IEC/IEEE 12207:2021} para detalhamento dos processos de ciclo de vida técnico.}
}

\sintesecapitulo{%
  \itemcaixa{A engenharia do CineDoom comprova que arquitetura de alta concorrência e processos de ciclo de vida rigorosos são indispensáveis para garantir a integridade de vendas de ingressos digitais.}
}
